\section{Temporal calibration and SI report maps}
\label{app:manual-temporal-calibration}

Counts at a declared cut level \(\ell\) are recorded in \(\mathrm{bip}_{\ell}\), and reciprocal rates at that level are recorded in \(\mathrm{biz}_{\ell}\). The manual may use the reference convention \(\ell=0\), written \(\mathrm{monon}_0\), \(\mathrm{bip}_0\), and \(\mathrm{biz}_0\). The hierarchy has no bottom cut, and \(\mathrm{bip}_0\) is not Planck time.

SI values are calibration or report targets. They do not generate the AO field relation. Distance, mass, charge, velocity, and phase appear only after a declared report map or temporal normalization.

A mass-rate anchor may be recorded as
\[
m \mapsto m_z := \frac{mc^2}{h},
\qquad
n := \log_3(m_z),
\qquad
\lambda := \frac{n}{4}.
\]
By default this is a report-layer anchor. A chapter may admit it as an exact working quantity only by declaring that status explicitly. Momentum-like quantities must be named by their AOD role, such as \(p_{\mathrm{AOD}}\), \(p_{\mathrm{loss}}\), or \(p_{\mathrm{prox}}\), and must be defined before use.

The generic report map has type
\[
\Psi_R:
\mathrm{ExactRow}_{\ell}(\mathcal F_K,\omega)
\longrightarrow
\mathrm{ReportQuantity}.
\]
